\documentclass[11pt,a4paper]{article}

% --- Encoding ---
\usepackage[utf8]{inputenc}
\usepackage[T1]{fontenc}
\usepackage[spanish,es-noshorthands]{babel}
\usepackage{lmodern}

% --- Layout ---
\usepackage[margin=2.5cm]{geometry}
\usepackage{parskip}
\setlength{\parindent}{0pt}
\setlength{\parskip}{0.5em plus 0.2em minus 0.1em}

% --- Tables and graphics ---
\usepackage{booktabs}
\usepackage{array}
\usepackage{enumitem}
\usepackage{xcolor}
\usepackage{graphicx}
\usepackage{tikz}
\usetikzlibrary{positioning,shapes.geometric,arrows.meta}
\usepackage{caption}

% --- Code listings ---
\usepackage{listings}
\definecolor{codegray}{rgb}{0.96,0.96,0.96}
\definecolor{codeborder}{rgb}{0.85,0.85,0.85}
\definecolor{codecomment}{rgb}{0.30,0.55,0.30}
\definecolor{codekeyword}{rgb}{0.10,0.30,0.70}
\definecolor{codestring}{rgb}{0.65,0.15,0.10}

\lstdefinestyle{cstyle}{
    language=C,
    backgroundcolor=\color{codegray},
    basicstyle=\ttfamily\footnotesize,
    breaklines=true,
    frame=single,
    rulecolor=\color{codeborder},
    showstringspaces=false,
    keywordstyle=\color{codekeyword}\bfseries,
    commentstyle=\color{codecomment}\itshape,
    stringstyle=\color{codestring},
    numbers=left,
    numberstyle=\tiny\color{gray},
    numbersep=8pt,
    aboveskip=1em,
    belowskip=1em,
}

\lstdefinestyle{shellstyle}{
    language=bash,
    backgroundcolor=\color{codegray},
    basicstyle=\ttfamily\small,
    breaklines=true,
    frame=single,
    rulecolor=\color{codeborder},
    showstringspaces=false,
    keywordstyle=\color{codekeyword}\bfseries,
    commentstyle=\color{codecomment}\itshape,
    stringstyle=\color{codestring},
    aboveskip=1em,
    belowskip=1em,
    columns=fullflexible,
}

\lstset{style=shellstyle}

% --- Definition box ---
\usepackage[most]{tcolorbox}
\newtcolorbox{definicion}[1][]{
    colback=blue!5!white,
    colframe=blue!50!black,
    fonttitle=\bfseries,
    title=Definición,
    sharp corners=southwest,
    rounded corners=northwest,
    arc=2pt,
    boxrule=0.6pt,
    left=8pt,right=8pt,top=4pt,bottom=4pt,
    #1
}
\newtcolorbox{idea}[1][]{
    colback=orange!5!white,
    colframe=orange!60!black,
    fonttitle=\bfseries,
    title=Idea intuitiva,
    sharp corners=southwest,
    rounded corners=northwest,
    arc=2pt,
    boxrule=0.6pt,
    left=8pt,right=8pt,top=4pt,bottom=4pt,
    #1
}

% --- Hyperlinks ---
\usepackage[colorlinks=true,
            linkcolor=blue!60!black,
            urlcolor=blue!50!black,
            citecolor=blue!60!black]{hyperref}

% --- Title ---
\title{
    \vspace{-2.5em}
    \textbf{Resumen teórico --- Sistemas Operativos} \\[0.3em]
    \large Capítulos 1 a 3 de las Notas del curso y capítulo 1 de Tannenbaum \& Bos
}
\author{
    Tomás Rodeghiero \\
    \small Sistemas Operativos --- Universidad Nacional de Río Cuarto
}
\date{2026}

\begin{document}

\maketitle

\begin{abstract}
\noindent
Este documento es un resumen teórico de los temas iniciales de la materia
Sistemas Operativos de la UNRC. Se basa en los capítulos 1, 2 y 3 de las
\emph{Notas del curso} (Marcelo Arroyo) y en el capítulo 1 del libro
\emph{Modern Operating Systems} (Tannenbaum \& Bos). El objetivo es servir
como material de estudio: las definiciones formales se acompañan de
explicaciones intuitivas, ejemplos y conexiones entre conceptos, de modo
que el lector pueda preparar tanto el examen teórico como el Práctico 1
con una comprensión sólida y articulada de los fundamentos.
\end{abstract}

\tableofcontents
\bigskip

% =====================================================================
\section{Introducción}
% =====================================================================

Una computadora moderna es un conjunto enorme de componentes
heterogéneos: uno o varios procesadores, memoria principal, discos,
periféricos de entrada y salida, controladores, buses, redes. Si cada
programador tuviera que tratar directamente con todos esos detalles, el
desarrollo de aplicaciones sería virtualmente imposible. El
\textbf{sistema operativo} (SO) resuelve ese problema: es la capa de
software que se interpone entre el hardware y los programas de usuario,
les ofrece abstracciones uniformes y administra el uso de los recursos
disponibles.

\begin{definicion}
Un \textbf{sistema operativo} es la pieza de software que abstrae el
hardware, administra los recursos del sistema (CPU, memoria,
dispositivos, archivos) y ofrece a las aplicaciones servicios uniformes
mediante \emph{llamadas al sistema}.
\end{definicion}

Las dos formas habituales de mirar al SO son complementarias:

\begin{itemize}[leftmargin=*]
    \item Como \textbf{máquina extendida}: oculta la fealdad del
        hardware (registros del controlador de disco, bits de estado de
        cada periférico, particularidades de cada arquitectura) y
        presenta a los programas un modelo limpio basado en archivos,
        procesos, sockets, etc. Tannenbaum lo resume diciendo que ``el
        SO convierte lo feo en lo bello''.
    \item Como \textbf{administrador de recursos}: la CPU, la memoria y
        los dispositivos son recursos limitados que múltiples procesos
        quieren usar al mismo tiempo. El SO multiplexa el acceso de
        forma ordenada, justa y protegida.
\end{itemize}

\begin{idea}
Si pensamos en un editor de video, un navegador y un cliente de correo
abiertos a la vez en una notebook, ninguno de los tres se entera del
otro. El SO logra esa ilusión multiplexando el procesador, partiendo la
memoria en zonas privadas y arbitrando el acceso al disco y a la red.
\end{idea}

% =====================================================================
\section{Estructura general de un sistema de cómputo}
% =====================================================================

\subsection{Capas de software y hardware}

Las Notas del curso presentan una estratificación típica:

\begin{enumerate}[leftmargin=*]
    \item \textbf{Hardware:} CPUs, RAM, dispositivos.
    \item \textbf{Kernel del SO:} corre en \emph{modo kernel}, tiene
        acceso total al hardware.
    \item \textbf{Bibliotecas y utilidades del sistema:} \texttt{libc},
        shells, compiladores, servicios (\emph{daemons}). Corren en
        \emph{modo usuario} pero son piezas del entorno operativo.
    \item \textbf{Aplicaciones de usuario:} programas que interactúan con
        el SO mediante la API que provee \texttt{libc} y otras
        bibliotecas.
\end{enumerate}

La frontera entre el modo usuario y el modo kernel es central para la
seguridad y la estabilidad del sistema. Solo el código del kernel puede
ejecutar instrucciones privilegiadas como deshabilitar interrupciones,
acceder a puertos de dispositivos o manipular las tablas de páginas.

\subsection{Hardware mínimo de interés}

\begin{itemize}[leftmargin=*]
    \item \textbf{CPU.} Ejecuta el ciclo
        \emph{fetch--decode--execute}. Posee registros generales,
        registros especiales (program counter, stack pointer,
        \emph{Program Status Word}), un conjunto de instrucciones
        propio (x86, ARM, RISC-V, \ldots) y \emph{modos de operación}.
        Las CPUs modernas usan \emph{pipelines} y diseños
        \emph{superscalares} que ejecutan varias instrucciones en
        paralelo, e incluso son \emph{multi-core}.
    \item \textbf{ROM.} Memoria no volátil con el firmware de arranque
        (BIOS o UEFI).
    \item \textbf{RAM.} Memoria principal volátil donde reside lo que la
        CPU está ejecutando.
    \item \textbf{Buses y puertos de E/S.} Conectan la CPU con
        dispositivos. Los \emph{controladores} (controllers) ofrecen
        registros que el kernel programa para iniciar transferencias.
    \item \textbf{DMA.} Transfiere datos entre dispositivo y memoria sin
        intervención de la CPU; al terminar, dispara una interrupción.
    \item \textbf{Jerarquía de memoria.} Registros, caché L1/L2, RAM,
        disco, todo con capacidades crecientes y velocidades
        decrecientes.
\end{itemize}

\begin{idea}
La idea fundamental para entender lo que viene es: la CPU pasa la mayor
parte del tiempo \emph{esperando}. Sin algún mecanismo que aproveche el
tiempo libre durante una operación de E/S, la máquina sería un
desperdicio. El SO se construye alrededor de esa observación.
\end{idea}

% =====================================================================
\section{Historia y evolución de los sistemas operativos}
% =====================================================================

Conviene mirar la evolución para entender por qué los conceptos modernos
existen.

\begin{enumerate}[leftmargin=*]
    \item \textbf{Primera generación (1945--1955), tubos de vacío.}
        No había SO. Los programas eran \emph{auto-contenidos}: cada
        uno incluía su propio código de E/S y rutinas matemáticas. Se
        cargaban manualmente vía \emph{plugboards} o cintas.
    \item \textbf{Segunda generación (1955--1965), transistores y batch
        systems.} Aparecen los \emph{lotes (batch)}: una pila de
        \emph{jobs} se ejecuta secuencialmente. El SO incipiente
        gestiona la sucesión de comandos
        (\$JOB, \$FORTRAN, \$LOAD, \$RUN, \$END).
    \item \textbf{Tercera generación (1965--1980), circuitos integrados
        y multiprogramación.} La CPU se vuelve más rápida que la E/S, así
        que se carga \emph{más de un programa} en memoria a la vez y, al
        bloquearse uno por E/S, otro toma la CPU. Esta es la idea de
        \textbf{multiprogramación}. El \emph{spooling} permite leer
        nuevos jobs desde disco. Aparece el \textbf{tiempo compartido
        (timesharing)} para dar respuesta interactiva
        (CTSS, MULTICS). Surge UNIX.
    \item \textbf{Cuarta generación (1980--presente), computadoras
        personales.} CP/M, MS-DOS, MacOS, Windows, Linux. Los SO se
        vuelven \emph{multiusuario} y \emph{multitarea} en hardware
        protegido. Se popularizan las interfaces gráficas (GUI).
    \item \textbf{Quinta generación (1990--presente), móviles.} Symbian,
        iOS, Android. Misma teoría, pero adaptada a recursos limitados,
        consumo y entornos táctiles.
\end{enumerate}

Una idea recurrente que aparece en el libro de Tannenbaum: las
tecnologías y técnicas se reinventan cíclicamente (``ontogeny
recapitulates phylogeny''). Cada nueva clase de dispositivos
(minicomputadoras, micros, smartphones, smartcards) suele empezar
\emph{primitiva} y va incorporando, a medida que el hardware lo
permite, las mismas técnicas que ya existían en mainframes.

% =====================================================================
\section{Funciones y objetivos del sistema operativo}
% =====================================================================

\subsection{Objetivos generales}

\begin{itemize}[leftmargin=*]
    \item \textbf{Abstraer el hardware.} Por ejemplo, en UNIX los
        dispositivos aparecen como \emph{archivos especiales} en
        \texttt{/dev}, y el almacenamiento masivo se modela como un
        \emph{sistema de archivos} con directorios y nombres
        simbólicos.
    \item \textbf{Administrar recursos.} Multiplexa la CPU entre
        procesos, particiona la memoria, arbitra el acceso a los
        dispositivos.
    \item \textbf{Proveer servicios.} Las \emph{llamadas al sistema}
        ofrecen E/S, manejo de procesos, archivos, comunicación, redes,
        sincronización.
    \item \textbf{Proteger el sistema.} Aislar los procesos entre sí y
        proteger al kernel de errores y abusos del software de usuario.
\end{itemize}

\subsection{Multiplexión en tiempo y en espacio}

Tannenbaum distingue dos modos de compartir un recurso:

\begin{itemize}[leftmargin=*]
    \item \textbf{Multiplexión en tiempo:} los usuarios o procesos se
        turnan. La CPU es un caso paradigmático: en un instante dado
        solo puede ejecutar un proceso, así que el SO conmuta entre
        ellos rápidamente. Otro ejemplo es la cola de impresión.
    \item \textbf{Multiplexión en espacio:} cada cliente recibe una
        \emph{porción} del recurso al mismo tiempo. La memoria principal
        se divide entre procesos y el espacio de disco entre archivos.
\end{itemize}

\subsection{Tipos de sistemas operativos}

Las Notas y el libro coinciden en un \emph{zoo} de variantes:

\begin{itemize}[leftmargin=*]
    \item \textbf{Mainframes:} priorizan el procesamiento de muchos jobs
        en paralelo (batch, transaccional, timesharing).
    \item \textbf{Servidores:} multiusuario, atienden clientes por red.
    \item \textbf{Multiprocesadores:} con varios CPUs (o cores) en una
        sola máquina.
    \item \textbf{PCs y portátiles (escritorio):} buscan dar buen
        soporte a un usuario interactivo.
    \item \textbf{Móviles:} multiplican GUIs táctiles, gestión agresiva
        de la batería y servicios de red.
    \item \textbf{Embedded y RTOS:} corren en dispositivos donde el
        software es fijo (microondas, automóviles, control industrial).
        Los \emph{RTOS} garantizan respuestas dentro de plazos
        acotados; pueden ser \emph{hard} (deadlines absolutos) o
        \emph{soft}.
    \item \textbf{Sensor-node y smartcards:} extremadamente reducidos.
\end{itemize}

% =====================================================================
\section{Kernel, modo usuario y modo kernel}
% =====================================================================

\subsection{El kernel}

\begin{definicion}
El \textbf{kernel} es el núcleo del sistema operativo: la porción de
código que corre en \emph{modo supervisor} y tiene acceso completo al
hardware. Incluye el manejo de procesos, la gestión de memoria, los
\emph{device drivers}, el sistema de archivos y la pila de red.
\end{definicion}

Un programa de usuario común \emph{no} es parte del kernel, aun si es
fundamental (la \texttt{libc}, el shell, los servicios de impresión
viven en \emph{user space}). La frontera precisa puede ser borrosa: en
sistemas tipo microkernel hasta el sistema de archivos puede vivir en
modo usuario.

\subsection{Modos de operación de la CPU}

Las CPUs modernas implementan al menos dos \emph{niveles de privilegio}:

\begin{itemize}[leftmargin=*]
    \item \textbf{Modo kernel (supervisor).} Toda instrucción es legal:
        habilitar/deshabilitar interrupciones, acceder a puertos,
        modificar registros de control de memoria, cambiar el modo.
    \item \textbf{Modo usuario.} Solo un subconjunto seguro de
        instrucciones está permitido. Cualquier intento de ejecutar una
        instrucción privilegiada genera una excepción y el kernel toma
        el control.
\end{itemize}

El bit que determina el modo está en el \emph{Program Status Word}
(PSW). El cambio de modo solo se produce mediante mecanismos
controlados: una interrupción, una excepción o una llamada al sistema.

\subsection{Mecanismos de protección}

Para que la multiprogramación sea segura se necesitan, según las Notas:

\begin{itemize}[leftmargin=*]
    \item \textbf{Confinamiento:} cada proceso solo puede acceder a su
        propio espacio de memoria. La CPU genera una \emph{segmentation
        fault} si lo viola.
    \item \textbf{Control:} los procesos no deben poder ejecutar
        instrucciones privilegiadas que les permitan tomar el control
        del kernel.
\end{itemize}

\begin{idea}
Sin protección, un \emph{bug} en un programa podría apagar el sistema o
robar datos de otro. El hardware moderno (MMU, anillos de privilegio)
está diseñado precisamente para ofrecerle al SO los mecanismos
necesarios para implementar esos invariantes.
\end{idea}

% =====================================================================
\section{Arquitecturas de kernel}
% =====================================================================

Tanto las Notas como el libro discuten varios estilos de diseño:

\begin{itemize}[leftmargin=*]
    \item \textbf{Monolíticos.} Todos los subsistemas se compilan y
        enlazan en un único ejecutable; se invocan entre sí como
        funciones ordinarias. Son simples y eficientes, pero un bug en
        un módulo puede afectar a todo el kernel. Ejemplos:
        UNIX clásico, Linux, FreeBSD.
    \item \textbf{Microkernels.} El kernel es muy pequeño y solo provee
        mecanismos básicos (planificación, IPC). Los subsistemas
        (filesystems, drivers, redes) corren como procesos
        independientes que se comunican por mensajes. Más robustos
        pero con sobrecarga de IPC. Ejemplos: Minix, L4, MS-Windows
        moderno (en parte).
    \item \textbf{Híbridos.} Combinan ambas ideas; ejemplo: macOS, que
        sobre el microkernel Mach corre buena parte de BSD UNIX como
        servicios.
    \item \textbf{Máquinas virtuales.} Un \emph{hypervisor} multiplexa
        el hardware entre varios SO huéspedes, cada uno corriendo
        encima de su propio entorno virtual. Ejemplos: KVM,
        Xen, VMware.
    \item \textbf{Contenedores.} \emph{No son virtualización}: comparten
        un mismo kernel, pero aíslan grupos de procesos y recursos
        (filesystems, redes virtuales, namespaces). En GNU-Linux se
        implementan con \emph{cgroups}; ejemplo de uso: Docker.
\end{itemize}

Un principio de diseño recurrente: separar \emph{mecanismos} (en las
capas bajas) de \emph{políticas} (en las capas altas) para favorecer la
flexibilidad. Por ejemplo, el kernel implementa el mecanismo de cambio
de contexto y el planificador define la política.

% =====================================================================
\section{Booting: arranque del sistema}
% =====================================================================

Resumen del proceso descripto en las Notas:

\begin{enumerate}[leftmargin=*]
    \item Al encender, la CPU comienza a ejecutar código en una
        dirección fija de la ROM: el \emph{firmware} (BIOS o UEFI).
    \item El firmware verifica RAM, dispositivos, y carga el primer
        sector del disco de arranque (MBR) o, en UEFI, un ejecutable
        de la partición EFI.
    \item Ese código suele ser un \emph{bootloader} (GRUB) que carga
        el kernel del SO en memoria y le transfiere el control.
    \item El kernel inicializa subsistemas (drivers, planificador,
        sistemas de archivos) y, finalmente, lanza el primer proceso de
        usuario (en UNIX: \texttt{init}, \texttt{systemd} o
        \texttt{launchd}). Ese proceso es el ancestro de todos los
        demás.
\end{enumerate}

% =====================================================================
\section{API, ABI y POSIX}
% =====================================================================

\begin{definicion}
La \textbf{API} (\emph{Application Programming Interface}) es la
interfaz a nivel de código fuente: declaraciones de funciones, tipos de
datos y constantes que el programador usa para invocar servicios del
sistema.
\end{definicion}

\begin{definicion}
La \textbf{ABI} (\emph{Application Binary Interface}) es la interfaz de
bajo nivel que especifica cómo se pasan argumentos, cómo se llaman las
funciones, qué registros se preservan, y qué formato tienen los
ejecutables (\texttt{ELF} en GNU-Linux, \texttt{Mach-O} en macOS,
\texttt{PE} en Windows).
\end{definicion}

\begin{definicion}
\textbf{POSIX} (\emph{Portable Operating System Interface}) es un
estándar del IEEE que define una API mínima común para los sistemas
tipo UNIX. Especifica syscalls (\texttt{open}, \texttt{read},
\texttt{write}, \texttt{fork}, \texttt{exec}, \ldots), utilidades de
shell y semánticas asociadas. Un programa que respeta POSIX corre, con
mínimas adaptaciones, en cualquier UNIX certificado.
\end{definicion}

% =====================================================================
\section{Programas, procesos y espacios de direcciones}
% =====================================================================

\subsection{Programa vs.\ proceso}

Un \textbf{programa} es un archivo ejecutable: una secuencia de
instrucciones de máquina y datos almacenados en disco. Un
\textbf{proceso}, en cambio, es ese programa \emph{en ejecución}, con
todo el contexto que necesita: registros, memoria, archivos abiertos,
etc.

\begin{definicion}
Un \textbf{proceso} es una instancia de un programa en ejecución. En un
sistema multitarea pueden coexistir múltiples instancias del mismo
programa.
\end{definicion}

\subsection{Imagen de memoria de un proceso}

En UNIX un proceso tiene tres áreas principales (segmentos):

\begin{itemize}[leftmargin=*]
    \item \textbf{Texto (text):} las instrucciones del programa. Es de
        solo lectura y se puede compartir entre instancias.
    \item \textbf{Datos (data):} variables globales y constantes
        inicializadas o no. Crece con \texttt{brk}/\texttt{sbrk} (y, en
        la práctica, vía \texttt{malloc}).
    \item \textbf{Pila (stack):} marcos de activación de funciones, con
        parámetros, variables locales y direcciones de retorno. Crece
        en sentido contrario a los datos.
\end{itemize}

Entre la pila y los datos queda un hueco de direcciones libres. La pila
crece hacia abajo (direcciones decrecientes) y los datos hacia arriba.

\subsection{Process control block (PCB)}

El kernel representa cada proceso con una estructura llamada \emph{PCB}
(o \emph{task descriptor}) que contiene:

\begin{itemize}[leftmargin=*]
    \item Identificador del proceso (\texttt{pid}).
    \item Estado: \texttt{RUNNING}, \texttt{READY}, \texttt{WAITING},
        \texttt{ZOMBIE}, \ldots
    \item Valor de los registros, incluido el \emph{program counter} y
        el \emph{stack pointer}.
    \item Lista de áreas de memoria en uso.
    \item Tabla de archivos abiertos (descriptores).
    \item Información de seguridad (UID, GID), señales pendientes,
        prioridad, padre, hijos.
\end{itemize}

\subsection{Espacio de direcciones y memoria virtual}

Cada proceso ve su propio espacio de direcciones lineal, típicamente
desde la dirección 0 hasta un máximo (por ejemplo, $2^{64}-1$ en una
plataforma de 64 bits). La \emph{MMU} (\emph{Memory Management Unit})
traduce esas direcciones \emph{virtuales} a direcciones \emph{físicas}
de RAM, posibilitando además la \textbf{memoria virtual}: páginas que
no se usan se llevan al disco y se traen a demanda. Esto permite
ejecutar programas más grandes que la memoria física.

\subsection{Árbol de procesos}

Cuando un proceso crea hijos, los hijos pueden a su vez crear nietos,
formando un árbol. Ejemplo típico: \texttt{init} es la raíz, lanza
\texttt{getty}/\texttt{login}, que lanzan el shell del usuario, que a
su vez lanza los comandos.

% =====================================================================
\section{Llamadas al sistema (syscalls)}
% =====================================================================

\subsection{Definición y mecanismo}

\begin{definicion}
Una \textbf{llamada al sistema} es una función que un proceso de
usuario invoca para pedir un servicio al kernel. Se implementan con
una instrucción especial (\emph{trap}, \emph{syscall},
\emph{software interrupt}) que provoca un cambio de modo.
\end{definicion}

El proceso típico de una syscall, según la figura 1-17 del libro:

\begin{enumerate}[leftmargin=*]
    \item El programa apila los parámetros en la pila o los pone en
        registros, según la ABI.
    \item Llama a la función de biblioteca correspondiente (por
        ejemplo, \texttt{read} de \texttt{libc}).
    \item Esa función coloca el \emph{número de syscall} en un
        registro convenido.
    \item Ejecuta la instrucción \texttt{TRAP} (o \texttt{syscall}, en
        x86-64).
    \item La CPU cambia a modo kernel y salta a un punto fijo del
        kernel: el \emph{trap handler}.
    \item El handler usa el número para indexar una tabla de
        despachadores y llama a la rutina del kernel que corresponde.
    \item Cuando la rutina termina, el kernel restaura el contexto y
        retorna al punto siguiente al \texttt{TRAP}, ya en modo
        usuario.
\end{enumerate}

\begin{idea}
Una syscall es \emph{como} una llamada a función ordinaria, pero con
dos diferencias importantes: (i) cambia de modo de privilegio y (ii) no
salta a una dirección arbitraria, sino a un punto fijo del kernel.
\end{idea}

\subsection{Las syscalls POSIX más usadas}

A continuación, una selección de las que aparecen en las Notas y en el
libro (Figura 1-18 de Tannenbaum), agrupadas por área:

\begin{table}[h]
\centering
\small
\begin{tabular}{p{4.5cm} p{9cm}}
\toprule
\textbf{Syscall} & \textbf{Función} \\
\midrule
\multicolumn{2}{l}{\emph{Procesos}} \\
\midrule
\texttt{pid = fork()} & Crea un proceso hijo, copia exacta del padre. \\
\texttt{exit(status)} & Finaliza el proceso con el código indicado. \\
\texttt{waitpid(pid,\&st,opts)} & Espera la terminación de un hijo. \\
\texttt{execve(path,argv,envv)} & Reemplaza la imagen de memoria por
otro programa. \\
\texttt{getpid() / getppid()} & Devuelve el PID propio o del padre. \\
\texttt{kill(pid, sig)} & Envía una señal a un proceso. \\
\texttt{sleep(n)} & Bloquea al proceso \texttt{n} segundos. \\
\midrule
\multicolumn{2}{l}{\emph{Archivos}} \\
\midrule
\texttt{fd = open(path,mode)} & Abre un archivo y devuelve un
descriptor. \\
\texttt{n = read(fd,buf,n)} & Lee \texttt{n} bytes en el buffer. \\
\texttt{n = write(fd,buf,n)} & Escribe \texttt{n} bytes desde el
buffer. \\
\texttt{close(fd)} & Cierra el descriptor. \\
\texttt{lseek(fd,off,whence)} & Mueve el offset interno de E/S. \\
\texttt{stat(path,\&st)} & Obtiene metadatos del archivo. \\
\texttt{dup(fd)} & Duplica un descriptor en el siguiente slot libre. \\
\texttt{pipe(p)} & Crea un pipe sin nombre. \\
\midrule
\multicolumn{2}{l}{\emph{Directorios y filesystems}} \\
\midrule
\texttt{chdir(dir)} & Cambia el directorio de trabajo. \\
\texttt{mkdir(name,mode)} & Crea un directorio. \\
\texttt{link(f1,f2)} & Crea un \emph{hard link}. \\
\texttt{unlink(name)} & Elimina una entrada de directorio. \\
\texttt{mount/umount} & Monta o desmonta un filesystem. \\
\midrule
\multicolumn{2}{l}{\emph{Misceláneas}} \\
\midrule
\texttt{chmod(name,mode)} & Cambia los bits de protección. \\
\texttt{time(\&t)} & Devuelve segundos desde el 1/1/1970. \\
\texttt{alarm(n)} & Programa la señal \texttt{SIGALRM}. \\
\bottomrule
\end{tabular}
\caption{Algunas de las syscalls POSIX más comunes.}
\end{table}

% =====================================================================
\section{Interrupciones y excepciones}
% =====================================================================

Las interrupciones son la forma estándar en que el hardware notifica
eventos asincrónicos al software:

\begin{itemize}[leftmargin=*]
    \item Un dispositivo termina una transferencia y dispara una
        \emph{interrupt request} (IRQ).
    \item El controlador de interrupciones le avisa a la CPU.
    \item La CPU termina la instrucción actual, salva el PC y el PSW
        en la pila, conmuta a modo kernel y salta al \emph{interrupt
        handler} indexado en el \emph{interrupt vector}.
    \item El handler atiende la interrupción y retorna; la CPU
        restaura el contexto y continúa el flujo previo.
\end{itemize}

Las \emph{excepciones} son disparadas por la propia CPU en respuesta a
condiciones del programa: división por cero, acceso a memoria
prohibida, instrucción ilegal. El kernel decide si terminar el
proceso, ignorar el problema o convertirlo en una señal manejable.

Un detalle importante: las syscalls también se implementan con un
mecanismo similar al de las interrupciones (instrucción \texttt{TRAP}
o \emph{software interrupt}), por lo que comparten infraestructura.

\subsection{Métodos de E/S}

Tannenbaum distingue tres formas de hacer entrada/salida:

\begin{itemize}[leftmargin=*]
    \item \textbf{Busy waiting (polling):} el driver inicia la
        operación y se queda en un bucle preguntando por el estado del
        dispositivo. Simple, pero desperdicia CPU.
    \item \textbf{Por interrupciones:} el driver inicia la operación y
        retorna; cuando el dispositivo termina, dispara una
        interrupción que despierta al proceso.
    \item \textbf{DMA:} el chip DMA copia datos entre dispositivo y
        memoria sin intervenir la CPU; al terminar, una interrupción
        cierra el ciclo.
\end{itemize}

% =====================================================================
\section{Multiprogramación, multitarea y tiempo compartido}
% =====================================================================

\begin{definicion}
La \textbf{multiprogramación} es la técnica de mantener varios
procesos cargados en memoria al mismo tiempo y multiplexar la CPU
entre ellos para mantenerla ocupada incluso durante operaciones de
E/S.
\end{definicion}

\begin{definicion}
La \textbf{multitarea} es la propiedad de un SO de ejecutar
aparentemente varios procesos en simultáneo. Puede ser
\emph{cooperativa} (los procesos ceden voluntariamente la CPU) o
\emph{preemptiva} (el kernel les puede sacar la CPU).
\end{definicion}

\begin{definicion}
El \textbf{tiempo compartido (timesharing)} es una variante de
multiprogramación interactiva en la que el kernel, ante una
interrupción periódica del reloj, puede quitarle la CPU al proceso
corriente y dársela a otro.
\end{definicion}

El cambio de proceso se llama \textbf{context switch}. Implica salvar
el estado del proceso saliente (registros, PC, PSW, mappings de la
MMU) en su PCB y restaurar el del proceso entrante. Es una operación
relativamente costosa que el SO trata de minimizar.

\begin{idea}
La sensación de simultaneidad surge porque el SO multiplexa muy
rápido. En un único instante, una CPU mononúcleo solo está ejecutando
\emph{un} proceso; la ilusión de paralelismo se construye sobre el
hecho de que cada proceso recibe muchos rebanadas pequeñas de tiempo
por segundo.
\end{idea}

% =====================================================================
\section{El shell y la sesión de usuario}
% =====================================================================

\subsection{Sesiones, login y multiusuario}

Un sistema multiusuario maneja \emph{sesiones}. Una terminal autentica
al usuario (\emph{login}), levanta un shell con sus credenciales
(\texttt{UID}, \texttt{GID}, \texttt{HOME}) y queda lista para recibir
comandos. El cierre de la sesión (\texttt{Ctrl-D} o \texttt{exit})
termina ese shell; \texttt{init} reinicia el ciclo en esa terminal.

\subsection{El shell}

\begin{definicion}
Un \textbf{shell} es un intérprete REPL (read--eval--print loop) que
recibe comandos del usuario, los analiza y ejecuta el programa
correspondiente. Implementa además operadores que combinan procesos
(redirección, secuencia, pipes, condicionales).
\end{definicion}

Cada vez que el usuario escribe un comando, el shell hace
\texttt{fork}, en el hijo aplica las redirecciones pedidas y ejecuta
\texttt{execve}; el padre espera con \texttt{waitpid}. El esqueleto
mínimo aparece en Tannenbaum (Figura 1-19) como:

\begin{lstlisting}[style=cstyle]
while (TRUE) {
    type_prompt();
    read_command(command, parameters);
    if (fork() != 0) {
        waitpid(-1, &status, 0);   /* padre espera al hijo */
    } else {
        execve(command, parameters, 0);  /* hijo ejecuta el comando */
    }
}
\end{lstlisting}

\subsection{Operadores del shell}

\begin{itemize}[leftmargin=*]
    \item \textbf{Secuenciales:} \texttt{cmd1 ; cmd2} ejecuta los dos
        en orden, sin importar el éxito.
    \item \textbf{Condicionales:} \texttt{cmd1 \&\& cmd2} ejecuta
        \texttt{cmd2} solo si \texttt{cmd1} terminó con
        \texttt{exit\_code = 0}; \texttt{cmd1 || cmd2} hace lo
        contrario.
    \item \textbf{Background:} \texttt{cmd \&} lanza el comando como
        proceso de fondo; el shell asigna un \emph{job id} y devuelve
        el control de inmediato.
    \item \textbf{Redirección:} \texttt{cmd > file} (sobreescribe),
        \texttt{cmd >> file} (append), \texttt{cmd < file}, \texttt{cmd
        2> file} (redirige stderr).
    \item \textbf{Pipes:} \texttt{cmd1 | cmd2} envía la salida estándar
        del primero a la entrada estándar del segundo.
\end{itemize}

\subsection{Variables de ambiente y \texttt{PATH}}

Cada sesión tiene un conjunto de \emph{variables de ambiente} que se
heredan a los procesos hijos en el tercer parámetro de \texttt{main}
(\texttt{char *envv[]}). Las relevantes:

\begin{itemize}[leftmargin=*]
    \item \texttt{HOME}: directorio personal.
    \item \texttt{PWD}: directorio actual.
    \item \texttt{USER}: usuario actual.
    \item \texttt{SHELL}: intérprete configurado.
    \item \texttt{PATH}: lista de directorios separados por
        \texttt{:} donde el shell busca ejecutables. Si el comando se
        invoca con barras (\texttt{./prog} o \texttt{/usr/bin/prog}) se
        toma como ruta directa; si no, se busca en \texttt{PATH}.
\end{itemize}

% =====================================================================
\section{Archivos y sistemas de archivos}
% =====================================================================

\subsection{Archivos como secuencia de bytes}

Para la mayoría de los SO modernos, un archivo es simplemente una
secuencia de bytes. La estructura interna (texto, ELF, JPEG, \ldots)
la imponen las aplicaciones, no el kernel. Los archivos se identifican
por su \emph{path}, que puede ser \emph{absoluto} (empieza en la raíz
\texttt{/}) o \emph{relativo} al directorio de trabajo.

Cada archivo tiene metadatos (tamaño, propietario, permisos,
timestamps, contador de enlaces, tipo) almacenados en una estructura
llamada \textbf{inodo}. La entrada de directorio asocia un nombre con
un inodo.

\subsection{Tipos de archivos en UNIX}

\begin{itemize}[leftmargin=*]
    \item \textbf{Regulares:} datos de usuario, código binario, texto.
    \item \textbf{Directorios:} contienen pares
        \emph{(nombre, inodo)}. Cada uno tiene siempre \texttt{.} y
        \texttt{..}.
    \item \textbf{Archivos especiales} (en \texttt{/dev}): representan
        dispositivos. Pueden ser \emph{block special} (discos) o
        \emph{character special} (terminales, mouse).
    \item \textbf{Pipes con nombre (FIFO):} archivos especiales para
        comunicación entre procesos no emparentados.
    \item \textbf{Enlaces:} ya sean \emph{symbolic} (apuntan por ruta)
        o \emph{hard} (otra entrada al mismo inodo).
\end{itemize}

\subsection{Descriptores de archivo}

\begin{definicion}
Un \textbf{descriptor de archivo} es un entero no negativo que el
kernel le asigna a un proceso al abrir un recurso. El proceso lo usa
en las syscalls posteriores (\texttt{read}, \texttt{write},
\texttt{close}, \texttt{lseek}).
\end{definicion}

Cada proceso, al ser lanzado por \texttt{init}, hereda tres
descriptores conectados a la terminal:

\begin{itemize}[leftmargin=*]
    \item \texttt{0}: \emph{stdin} (entrada estándar, teclado).
    \item \texttt{1}: \emph{stdout} (salida estándar, pantalla).
    \item \texttt{2}: \emph{stderr} (salida de errores, pantalla).
\end{itemize}

La syscall \texttt{dup(fd)} duplica un descriptor en el primer slot
libre; combinada con \texttt{close} y \texttt{open}, es la base de las
redirecciones de E/S del shell.

\subsection{Filesystem jerárquico}

Las Notas y el libro coinciden en describir el filesystem como un
árbol con un único directorio raíz \texttt{/}. Los filesystems
removibles (CDs, USB, particiones) se incorporan al árbol mediante la
syscall \texttt{mount}, que los \emph{ata} a un directorio existente.
\texttt{umount} hace lo contrario.

\subsection{Permisos UNIX}

Cada archivo tiene un \emph{modo} de 9 bits agrupados en tres
ternas (dueño, grupo, otros) con los flags \texttt{r} (lectura),
\texttt{w} (escritura) y \texttt{x} (ejecución). Por ejemplo, el modo
\texttt{rwxr--r--} permite al dueño leer, escribir y ejecutar; al
grupo y a otros, solo leer. Para directorios, \texttt{x} significa
permiso para \emph{atravesarlos}.

% =====================================================================
\section{Pipes y FIFOs (IPC)}
% =====================================================================

\subsection{Pipes sin nombre}

\begin{definicion}
Un \textbf{pipe} es un buffer acotado en el kernel con dos extremos
(lectura y escritura) que implementa el patrón productor/consumidor.
La syscall \texttt{pipe(p)} crea uno y devuelve dos descriptores:
\texttt{p[0]} (lectura) y \texttt{p[1]} (escritura).
\end{definicion}

Los pipes solo permiten comunicar procesos \emph{emparentados} porque
los descriptores se comparten vía \texttt{fork}. El comando shell
\texttt{cmd1 | cmd2} se implementa así:

\begin{lstlisting}[style=cstyle]
int p[2];
pipe(p);
if (fork() == 0) {                 /* primer hijo: cmd1 */
    close(p[0]);                   /* no lee del pipe */
    dup2(p[1], STDOUT_FILENO);     /* su stdout va al pipe */
    close(p[1]);
    execvp("cmd1", args1);
}
if (fork() == 0) {                 /* segundo hijo: cmd2 */
    close(p[1]);                   /* no escribe en el pipe */
    dup2(p[0], STDIN_FILENO);      /* su stdin viene del pipe */
    close(p[0]);
    execvp("cmd2", args2);
}
close(p[0]); close(p[1]);
wait(NULL); wait(NULL);
\end{lstlisting}

\textbf{Importante:} si el padre no cierra ambos extremos, el lector
nunca verá EOF y se quedará bloqueado para siempre.

\subsection{FIFOs (pipes con nombre)}

Un \textbf{FIFO} es un pipe que tiene un nombre en el sistema de
archivos. Se crea con \texttt{mkfifo(path, mode)} y se abre con
\texttt{open}/\texttt{read}/\texttt{write}/\texttt{close}. Permite
comunicar procesos \emph{independientes} (sin parentesco). Los datos
no se almacenan en disco; el archivo solo da un punto de encuentro y
los datos viven en el buffer del kernel mientras hay procesos
conectados.

% =====================================================================
\section{Señales y alarmas}
% =====================================================================

\subsection{Señales}

\begin{definicion}
Una \textbf{señal} es una notificación asíncrona de un evento enviada
por el kernel (o por otro proceso vía \texttt{kill}) a un proceso
destinatario. Es el ``equivalente software'' de una interrupción.
\end{definicion}

Algunos ejemplos:

\begin{itemize}[leftmargin=*]
    \item \texttt{SIGINT} (\texttt{Ctrl-C}): interrumpir el proceso.
    \item \texttt{SIGTSTP} (\texttt{Ctrl-Z}): suspenderlo.
    \item \texttt{SIGSEGV}: acceso a memoria inválida.
    \item \texttt{SIGFPE}: error aritmético (división por cero).
    \item \texttt{SIGALRM}: alarma vencida.
    \item \texttt{SIGKILL}: terminación incondicional. No se puede
        manejar ni ignorar.
\end{itemize}

Un proceso puede instalar un \emph{handler} con \texttt{signal} o
\texttt{sigaction}; cuando el kernel le entrega la señal, le
\emph{salta} al handler antes de continuar el flujo normal. Si no hay
handler, el kernel ejecuta la acción por defecto (terminar, ignorar,
suspender, según la señal).

\textbf{Restricción importante:} dentro de un handler solo deben
invocarse funciones \emph{async-signal-safe}: \texttt{read},
\texttt{write}, \texttt{\_exit}, \texttt{sigaction}, etc. Las
funciones de \texttt{stdio} (\texttt{printf}, \texttt{fgets},
\texttt{scanf}) usan buffers internos que pueden quedar en estado
inconsistente al ser interrumpidos.

\subsection{Alarmas}

La syscall \texttt{alarm(n)} le pide al kernel que envíe
\texttt{SIGALRM} luego de \texttt{n} segundos. Es la base para
implementar \emph{timeouts} sobre operaciones bloqueantes:

\begin{lstlisting}[style=cstyle]
signal(SIGALRM, on_alarm);
alarm(5);                         /* 5 segundos */
nread = read(STDIN_FILENO, buf, n);
alarm(0);                         /* cancelar si llegó dato */
\end{lstlisting}

% =====================================================================
\section{Hilos y concurrencia (introducción)}
% =====================================================================

Un proceso clásico tiene un único hilo de ejecución (un \emph{program
counter}, una pila). Un \textbf{thread} es un hilo de ejecución
independiente que comparte el espacio de direcciones con otros hilos
del mismo proceso. Los threads se introducen en el capítulo 1 como
``procesos livianos'' y permiten paralelizar dentro de una misma
aplicación. La sincronización (mutexes, semáforos, condiciones) es
necesaria para acceder a datos compartidos sin condiciones de carrera;
ese tema se desarrolla en prácticos posteriores.

% =====================================================================
\section{El lenguaje C y la cadena de herramientas}
% =====================================================================

Las Notas dedican el capítulo 3 a presentar C y la \emph{toolchain}
porque son el vehículo natural para programar el SO y las
aplicaciones que ejercitan sus syscalls. Conceptos relevantes:

\subsection{Tipos y construcciones}

\begin{itemize}[leftmargin=*]
    \item Tipos básicos: \texttt{char}, \texttt{int}, \texttt{long}
        (con/sin signo), \texttt{float}, \texttt{double}.
    \item Compuestos: arreglos, strings (terminados en \texttt{\textbackslash 0}),
        \texttt{struct}, \texttt{union}.
    \item Punteros y aritmética de punteros.
    \item Funciones: pasaje por valor, simulación de pasaje por
        referencia con punteros, funciones variádicas como
        \texttt{printf}.
    \item Módulos: separación en archivos \texttt{.c} y cabeceras
        \texttt{.h}; \texttt{\#include} y declaraciones \texttt{extern}.
\end{itemize}

\subsection{Pasos de compilación}

\begin{enumerate}[leftmargin=*]
    \item \textbf{Pre-procesado} (\texttt{cpp}): expande
        \texttt{\#include}, \texttt{\#define} y demás directivas.
    \item \textbf{Compilación} (\texttt{cc1}): traduce de C a
        \emph{assembly}.
    \item \textbf{Ensamblado} (\texttt{as}): genera archivos objeto
        \texttt{.o}.
    \item \textbf{Linking} (\texttt{ld}): combina objetos y bibliotecas
        en un ejecutable, resolviendo símbolos y referencias externas.
\end{enumerate}

\subsection{Bibliotecas estáticas y dinámicas}

\begin{itemize}[leftmargin=*]
    \item \textbf{Estáticas} (\texttt{libfoo.a}): el código de la
        biblioteca se incrusta en el ejecutable durante el linking.
        Resultado: ejecutables más grandes pero auto-contenidos.
    \item \textbf{Dinámicas} (\texttt{libfoo.so} en Linux,
        \texttt{.dylib} en macOS, \texttt{.dll} en Windows): el
        ejecutable solo guarda referencias; el \emph{dynamic linker}
        las resuelve en tiempo de carga (o, con \emph{lazy linking}, en
        la primera invocación). Ejecutables más chicos y biblioteca
        compartida en memoria entre procesos.
\end{itemize}

\subsection{\texttt{make} y los Makefiles}

\texttt{make} automatiza la compilación a partir de reglas
\texttt{target: dependencias}. Recompila solo lo necesario comparando
\emph{timestamps}. Las \emph{autotools} (\texttt{autoconf},
\texttt{automake}) generan Makefiles portables.

% =====================================================================
\section{Resumen integrador}
% =====================================================================

Para fijar las ideas, conviene visualizar cómo todos estos conceptos
encajan al ejecutar un comando de shell, por ejemplo:

\begin{lstlisting}
$ ls -l | wc -l > cuentas.txt
\end{lstlisting}

\begin{enumerate}[leftmargin=*]
    \item El usuario tipea el comando en la terminal. El \emph{tty}
        driver del kernel envía los caracteres al shell.
    \item El shell parsea la línea: detecta el pipe (\texttt{|}) y la
        redirección (\texttt{>}).
    \item Llama a \texttt{pipe(p)} para crear un pipe entre los dos
        comandos.
    \item Hace \texttt{fork} dos veces:
        \begin{itemize}
            \item El primer hijo redirige su \texttt{stdout} al
                extremo de escritura del pipe (\texttt{dup2}) y ejecuta
                \texttt{ls -l} con \texttt{execve}.
            \item El segundo hijo redirige su \texttt{stdin} al
                extremo de lectura, abre \texttt{cuentas.txt} con
                \texttt{O\_WRONLY|O\_CREAT|O\_TRUNC}, redirige su
                \texttt{stdout} a ese archivo y ejecuta \texttt{wc -l}.
        \end{itemize}
    \item El padre cierra ambos extremos del pipe y llama a
        \texttt{waitpid} dos veces para esperar a los hijos.
    \item Cada \texttt{execve} hace que el kernel cargue el ELF
        correspondiente, prepare el espacio de direcciones (texto,
        datos, pila), enlace dinámicamente \texttt{libc} y salte a
        \texttt{main}.
    \item Los procesos ejecutan, hacen \texttt{read} y \texttt{write}
        contra los descriptores que el shell les preparó. Cuando
        \texttt{ls} termina, cierra sus descriptores y el extremo de
        escritura del pipe se libera; \texttt{wc} ve EOF, termina y
        cierra su archivo de salida.
    \item El shell recoge los exit status, actualiza \texttt{\$?} y
        vuelve a mostrar el prompt.
\end{enumerate}

Detrás de un comando aparentemente simple operan, simultáneamente:
\emph{procesos}, \emph{multiplexión de CPU}, \emph{descriptores y
redirección}, \emph{pipes como IPC}, \emph{permisos del filesystem},
\emph{libc dinámica}, \emph{cambios de modo} para cada syscall y
\emph{interrupciones} cada vez que el kernel hace context switch o
atiende una E/S. El estudio de los sistemas operativos consiste,
precisamente, en entender cómo todas esas piezas funcionan en conjunto.

% =====================================================================
\section*{Referencias}
\addcontentsline{toc}{section}{Referencias}
% =====================================================================

\begin{enumerate}[leftmargin=*]
    \item Marcelo Arroyo. \emph{Notas del curso de Sistemas
        Operativos}, capítulos 1, 2 y 3. UNRC, 2026.\\
        \url{https://marceloarroyo.gitlab.io/cursos/SO/}
    \item Andrew S. Tanenbaum, Herbert Bos. \emph{Modern Operating
        Systems}, capítulo 1, 5\textsuperscript{a} edición. Pearson,
        2024.
    \item Remzi y Andrea Arpaci-Dusseau. \emph{Operating Systems: Three
        Easy Pieces}. \url{https://pages.cs.wisc.edu/~remzi/OSTEP/}
    \item Avi Silberschatz, Peter B. Galvin, Greg Gagne. \emph{Operating
        System Concepts}, 9\textsuperscript{a} edición.
    \item POSIX.1-2017 / Single UNIX Specification. The Open Group.
\end{enumerate}

\end{document}
